\section{Conclusion}

\lesson{Probably no unit at your institution ``owns'' exam
  administration: different pieces are owned by different units.}

Delivering exams securely and efficiently is not a simple function of
just the LMS, or just the firewall, or just the computer lab facilities, or just the
scheduling, or just the proctors; it requires coordination among all
these elements.
At most campuses, responsibility for any
one of them usually resides in a specific unit, but
nobody ``owns'' entire-course proctoring as a whole.
Achieving it requires new technical and social cooperation across
campus unit boundaries that are usually non-porous.   All of the
challenges we faced inherit from this challenge in some way.

Scaling up from our pilot has not been without additional challenges, including
getting the attention of the campus to recognize the value of
centrally funding a larger effort;
but it has led
to an assessment-delivery model that is (as others before us have found) 
\emph{far\/} superior to traditional exams in most respects.
The proof of the pudding is that the CBTF is now generating its own demand.

Our experience
demonstrates that the model pioneered at UIUC can be 
incrementally bootstrapped in an environment with different cultural
and administrative constraints without ``boiling the ocean.''  We are
confident that other institutions can do the same, and we
hope our lessons and reflections will help them do so.
